\section{Integration}
\thispagestyle{plain}

\subsection{Monte Carlo Integration}
Classical numerical methods for integration (e.g. Simpson's rule) become infeasible for high-dimensional integrals,
essentially, as the errors scale with the number of function-evaluation-points per dimension (so in high dimensions, we need
to perform lots of function evaluations).

\subsubsection{Intuition for Monte Carlo Integration}
Consider we want to calculate the integral

\begin{equation}
    I := \int_{V} g(\vec{x}) d^d \vec{x}
\end{equation}

where $V$ is a $d$-dimensional volume and $g$ is a function $g: V \rightarrow \mathbb{R}$. Then it makes intuitive sense
to approximate this integral by the mean of the function $g$ on the domain of interest, as found as the mean of $g$ evaluated
at $N$ random points $\vec{x}_i$ in $V$, multiplied by the volume of $V$ (also denoted by $V$)

\begin{equation}
\hat{I}_N := \frac{V}{N} \sum_{i=1}^{N} g(\vec{x}_i) \xrightarrow[N \rightarrow \infty]{a.s.} I
\end{equation}

\bluebox{\textbf{Monte Carlo Integration}: Let $V$ for simplicity be the hypercube $[0,1]^d$ (by change of variables, the domain of interest can be mapped to this hypercube). Then Monte Carlo
integration to approximate $I = \int_{[0,1]^d} g(\vec{x}) d^d \vec{x}, \vec{x} \in \mathbb{R}^d$ is given by
\begin{itemize}
    \item Generate $N$ random vectors $\vec{x}$ where each component is drawn uniformly from $[0,1]$ (so we need $N\times d$ random numbers)
    \item Evaluate
    \begin{equation}
        \hat{I}_N = V \frac{1}{N} \sum_{i=1}^{N} g(\vec{x}_i), \quad \text{here } V = 1 
    \end{equation}
\end{itemize}}

\greenbox{\textbf{Scaling of the error}: The error of the Monte Carlo integrator scales with $1/\sqrt{N}$, independent of the number of dimensions of the integration $d$.}

\subsubsection{Connection to the Monte Carlo Estimator}
We can also obtain the result from above
from the Monte Carlo Estimator (eq. \ref{eq:mc_est}) by using the uniform distribution
\begin{equation}
    f(\vec{x}) = \begin{cases}
        \frac{1}{V} & \text{if } \vec{x} \in V \\
        0 & \text{else}
    \end{cases}
\end{equation}
and rewriting the integral as
\begin{equation}
    I = \int_{V} g(\vec{x}) d^d \vec{x} = V \int_{V} g(\vec{x}) \cdot \frac{1}{V} d^d \vec{x} = V \int_{V} g(\vec{x}) \cdot f(\vec{x}) d^d \vec{x}
\end{equation}

\subsubsection{Intuition from calculating the area of a shape}
Consider an irregular shape with area $A_{\text{shape}}$ within a rectangle with area $A_{\text{rect}}$. 
The probability that a point drawn uniformly from the rectangle lies within the shape is 
\begin{equation}
    p_{\text{shape}} = \frac{A_{\text{shape}}}{A_{\text{rect}}}
\end{equation}
so we can approximate
\begin{equation}
    A_{\text{shape}}\approxeq A_{\text{rect}} \frac{\# \text{ hits in shape}}{\# \text{ hits in rectangle}}
\end{equation}
which is the same as the Monte Carlo Integration introduced, where $g(\vec{x}) = 1_{\text{shape}}(\vec{x})$ and $V = A_{\text{rect}}$.

\subsubsection{Error in Monte Carlo Integration - scaling with $\frac{1}{\sqrt{N}}$}
Based on the connection to the Monte Carlo Estimator, the
variance of our estimator for the integral is given by
\begin{equation}
    \Var[\hat{I}_N] = \frac{V^2}{N} \Var[g(x)], \quad \text{estimate} \Var[g(x)] \text{ by } S_{g;N}^2 := \frac{1}{N-1} \sum_{i=1}^{N} \left(g(x_i) - \frac{\hat{I}_N}{V}\right)^2
    \label{eq:mcint_var}
\end{equation}
\bluebox{
So the standard error of our Monte Carlo estimator
\begin{equation}
    \hat{I}_N = \frac{V}{N} \sum_{i=1}^{N} g(\vec{x}_i)
\end{equation}
that is if we calculate $\hat{I}_N$ multiple times with different random samples, the standard deviation of these
different estimates is
 \begin{equation}
    \boxed{\sigma_{\hat{I}_N} = V \sqrt{\frac{\sigma_g^2}{N}} \propto \frac{1}{\sqrt{N}}}
 \end{equation}
which is illustrated in Fig. \ref{fig:mcint}. This follows directly from the basic properties of the variance
\begin{equation}
    \text{for } X,Y \text{ independent} \quad \Var[X+Y] = \Var[X] + \Var[Y], \quad \Var[aX] = a^2 \Var[X]
\end{equation}
with
\begin{equation}
    \Var[\hat{I}_N] = \Var\left[\frac{V}{N} \sum_{i=1}^{N} g(\vec{x}_i)\right] = \frac{V^2}{N^2} \sum_{i=1}^{N} \Var[g(\vec{x}_i)] \underset{\vec{x}_i \text{ i.i.d.}}{=} \frac{V^2}{N} \Var[g(\vec{x})]
\end{equation}
}

\subsubsection{Distribution of $\hat{I}_N$ and derivation of the central limit theorem\skipthis}
The Monte Carlo Estimator is
\begin{equation}
    \hat{I}_N = \frac{V}{N} \sum_{i=1}^{N} g(\vec{x}_i) = \frac{1}{N} \sum_{i=1}^{N} y_i = \sum_{i=1}^{N} s_i, \quad y_i = V g(\vec{x}_i), \quad s_i = \frac{y_i}{N}
\end{equation}
by basic properties of variance and mean
\begin{equation}
    \begin{gathered}
        \langle \hat{I}_N \rangle = N \langle s \rangle = \langle y \rangle  \\
        \Var[\hat{I}_N] = N \Var[s] = \frac{1}{N} \Var[y] = \frac{V^2}{N} \Var[g(\vec{x})]
    \end{gathered}
\end{equation}
\note{$p(y)$ is the distribution of $y$ over $V$ along the $y$-axis with
\begin{equation}
    \langle y \rangle = \int y p(y) \, dy = \int_{V} g(\vec{x}) \, d^d \vec{x}, \quad \int p(y) \, dy = 1
\end{equation}
\textbf{$p(y)$ might be any distribution.}
}
Applying the Central Limit Theorem (eq. \ref{eq:central_limit_theorem}) to the Monte Carlo Estimator
(assuming the moments of $p(y)$ exist and are finite), we obtain
\begin{equation}
    \begin{aligned}
        P_N(\hat{I}_N) &= \frac{1}{\sqrt{2\pi \Var[\hat{I}_N]}} \exp{\left( -\frac{(\hat{I}_N-\langle \hat{I}_N \rangle)^2}{2\Var[\hat{I}_N]} \right)} \\
                       &= \frac{\sqrt{N}}{\sqrt{2\pi \sigma_y^2}} \exp{\left( -\frac{N(\hat{I}_N-\langle y \rangle)^2}{2\sigma_y^2} \right)}
    \end{aligned}
\end{equation}
independent of the shape of $p(y)$ for $N$ sufficiently large.

\paragraph{Derivation of the Central Limit Theorem by Fourier Transform} The central limit
theorem can be derived by considering $P_N$ in Fourier space, where 
the central ingredient will simply be
\begin{equation}
    \exp{x} = \lim_{N \rightarrow \infty} \left(1 + \frac{x}{N}\right)^N
\end{equation}
which is also approximately true for large $N$.
\bluebox{\textbf{Distribution of the sum of independent random variables}: Consider \textbf{independent} variables $X,Y$ with PDFs $f_X$ and $f_Y$. The PDF of the sum $Z = X + Y$
is naturally the convolution
\begin{equation}
    f_Z(z)=\int_{-\infty}^{\infty} \underbrace{f_X(x) f_Y(z-x)}_{x+(z-x) = z} d x
\end{equation}
Another way to write this is
\begin{equation}
    f_Z(z)=\int f_X(x) f_Y(y) \delta(z-(x+y)) \, d x \, d y
\end{equation}
which one can easily see is equivalent to the convolution and also makes intuitive sense - the $\delta$-function
ensures that the sum of $x$ and $y$ is $z$.
}
Analogously to the two-variable case, we write
\begin{equation}
    P_N\left(\hat{I}_N\right)=\int \delta\left(\hat{I}_N-\sum_i \frac{y_i}{N}\right) p\left(y_1\right) p\left(y_2\right) \cdots p\left(y_N\right) \mathrm{d} y_1 \mathrm{~d} y_1 \cdots \mathrm{d} y_N
\end{equation}
Based on the Fourier transform and expansion of $p(y)$
\begin{equation}
    \begin{aligned}
        \hat{p}(k)&=\int p(y) \exp{\left(i k(y-\langle y\rangle)\right)} \, \mathrm{d} y \\
        \underset{\text{Series Expansion}}&= \int p(y) \left[ 1 + ik(y-\langle y\rangle) - \frac{k^2}{2}(y-\langle y\rangle)^2 + \cdots \right] \, \mathrm{d} y \\
        &= 1 - \frac{k^2}{2} \sigma_y^2 + \cdots
    \end{aligned}
\end{equation}
we find the Fourier transform of $P_N$ to be
\begin{equation}
    \begin{aligned}
    \hat{P}_N(k) & =\int P_N\left(\hat{I}_N\right) \exp{\left(i k\left(I_N-\left\langle I_N\right\rangle\right)\right)} \mathrm{d} I_N \\
    \underset{\langle \hat{I}_N \rangle = \langle y \rangle}&= \int p\left(y_1\right) \cdots p\left(y_N\right) \exp{\left(i \frac{k}{N}\left(y_1-\left\langle y_1\right\rangle+y_2-\left\langle y_2\right\rangle+\ldots+y_N-\left\langle y_N\right\rangle\right)\right)} \mathrm{d} y_1 \cdots \mathrm{d} y_N \\
    &= \left[\hat{p}\left(\frac{k}{N}\right)\right]^N \\
    \underset{\text{series expansion of } \hat{p}}&= \left( 1 - \frac{k^2 \sigma_y^2}{2N^2} \right)^N \\
    \underset{N \text{ large}}&\approxeq \exp{\left( - \frac{k^2 \sigma_y^2}{2N} \right)}
    \end{aligned}
\end{equation}
where the inverse Fourier transform of $\hat{P}_N(k)$
\begin{equation}
    P_N\left(\hat{I}_N\right)=\frac{1}{2 \pi} \int \exp{\left(-i k\left(I_N-\left\langle I_N\right\rangle\right)\right)} \hat{P}_N(k) \mathrm{d} k
\end{equation}
yields the previously stated result for $P_N\left(\hat{I}_N\right)$\footnote{Insert the result for $\hat{P}_N(k)$, complete the square and use the normalization of the normal distribution
or the Gaussian integral $\int \exp \left(-\alpha x^2\right) \mathrm{d} x=\sqrt{\frac{\pi}{\alpha}}$.}.

\subsubsection{Illustrative Example of Monte Carlo Integration}
Consider the integral
\begin{equation}
    I = \int_{0}^{1} g(x) dx, \quad g(x) = \exp{\left(-\frac{x^2}{2}\right)}
\end{equation}
Our Monte Carlo estimator for this integral is given by
\begin{equation}
    \hat{I}_N = \frac{1}{N} \sum_{i=1}^{N} g(x_i), \quad x_i \sim \mathcal{U}(0, 1)
\end{equation}
which is illustrated in Fig. \ref{fig:mcint}. The higher the number of sampled points $N$ the lower the variance of our estimator; the lower the variance of $g$ along the $y$ axis, the lower the variance of our estimator.

\begin{figure}[!htb]
 \centering
 \includesvg[width=1\textwidth]{figures/mc_int.svg}\hfill
 \caption{Illustration of Monte Carlo Integration for $I = \int_{0}^{1} \exp{\left(-\frac{x^2}{2}\right)} dx$}
 \label{fig:mcint}
\end{figure}

\subsubsection{Comparison to other techniques - when to use Monte Carlo integration?}
\paragraph*{Error of standard methods} In a standard integration scheme like the midpoint rule
or Simpsons rule\footnote{$\int_a^b f(x) d x \approx \frac{b-a}{6}\left[f(a)+4 f\left(\frac{a+b}{2}\right)+f(b)\right]$
which can be derived by approximating the function $f$ by a quadratic polynomial and integrating this polynomial.} each
dimension is divided into $n$ regularly space points, so that we have $N = n^d$ points in total. The error scales with
\begin{equation}
    \text{midpoint or trapezoidal rule } \propto \frac{1}{n^2}, \quad \text{Simpson's rule } \propto \frac{1}{n^4} = \frac{1}{N^{\frac{4}{d}}}
\end{equation}
\redbox{In these \textit{regular} methods, adding more points to the integration (increasing $N$) helps less and less the higher we go 
in dimension - e.g. for Simpson's rule $\text{err} \propto \frac{1}{N^{\frac{4}{d}}}$. For a standard method if we want
$10$ points per dimension with $d = 10$ we already need $10^{10}$ function evaluations - infeasible (e.g. high-dimensional integrations in Machine Learning).}
\greenbox{On the other hand in Monte Carlo integration the error scales with $\frac{1}{N^\frac{1}{2}}$ independent of the dimensionality $d$.}
\bluebox{\textbf{Starting at what dimension of the integral $d$ will MC integration have a more
favorable error scaling compared to Simpson's rule?} Setting $\frac{1}{N^{\frac{4}{d}}} = \frac{1}{N^\frac{1}{2}}$ yields that starting at $d = 8$ 
Monte Carlo Integration's error will reduce more quickly if we increase $N$ the number of points (/ function evaluations) used for the integration.}

\paragraph{Intuition for the better scaling} One intuition for this advantage of Monte Carlo is that the
higher the dimension, the more evenly pairs of random points are spread with respect to their pair-wise distance
- normally known as the curse of dimensionality, illustrated in Fig. \ref{fig:curse}.

\begin{figure}[!htb]
 \centering
 \includesvg[width=0.8\textwidth]{figures/curse.svg}\hfill
 \caption{In higher dimensions, pairwise distances between random points show a sharper distribution around the mean distance
 (a more even distribution in space)}
 \label{fig:curse}
\end{figure}

\subsubsection{Reducing the variance of the Monte Carlo Estimator}

\subsubsubsection{Antithetic Estimators\skipthis}

Two random variables with the same distribution on the same probability space are called \textit{antithetic}, if their covariance is negative.
We can use such antithetic variables to reduce the variance of a Monte Carlo estimator.

This motivates the following ansatz

\begin{equation}
    \begin{multlined}
        I := \int_{0}^{1} g(x) dx, \quad \hat{I}_{anti} = \frac{\hat{I}_{N / 2} + \tilde{I}_{N / 2}}{2} \\
        \hat{I}_{N / 2} = \frac{1}{N / 2} \sum_{i=1}^{N / 2} g(u_i), \quad \tilde{I}_{N / 2} = \frac{1}{N / 2} \sum_{i=1}^{N / 2} g(1 - u_i), \quad u_i \sim \mathcal{U}(0, 1)
    \end{multlined}
\end{equation}

with

\begin{equation}
    \Var[\hat{I}_{anti}] = \frac{1}{4} \left(\Var[\hat{I}_{N / 2}] + \Var[\tilde{I}_{N / 2}] + 2 \Cov[\hat{I}_{N / 2}, \tilde{I}_{N / 2}]\right) \explain{=}{\text{see eq. \ref{eq:mcint_var}}} \Var[\hat{I}_N] + \frac{1}{2} \Cov[\hat{I}_{N / 2}, \tilde{I}_{N / 2}]
\end{equation}

where for $g$ continuous and monotonic, $\Cov[\hat{I}_{N / 2}, \tilde{I}_{N / 2}]$ is negative, as $U$ and $1 - U$ with $U \sim \mathcal{U}(0, 1)$ are negatively correlated. So we can reduce the variance compared to an estimator $\hat{I}_N$ using the same number of function 
evaluation and only half the number of random numbers.

Let us give an intuition why using antithetic evaluation points makes sense for estimating an integral of a monotonic function from $0$ to $1$. Using uniform antithetic variables we create pairs of points $(x_i, 1 - x_i)$, which are symmetric around $x = 0.5$ and as of the monotony of the function $g$ we want to integrate, we obtain a better balance between sampling points where $g$ is large and where $g$ is small, reducing the variance of the estimate.

\textbf{Example of using Antithetic Estimators}
\par
% Notice the even distribution of points around (a+b) / 2 by constructing yielding an advantage over standard M
Consider the integral
\begin{equation}
    I = \int_{0}^{2} \frac{1}{1 + x^2} dx \explain{=}{\text{exact solution}} \arctan(2) \approx 1.107149
\end{equation}
Estimates based on the standard Monte Carlo estimator and the antithetic estimator are given in table \ref{tab:antithetic} and illustrated in Fig. \ref{fig:mc_antithetic}. A relative reduction in standard deviation of roughly 80\% is achieved using the antithetic estimator.

% make table for results with N = 100 and seed = 1234 for results
% MC: est, std, conf. int: 1.09562114 0.01677497 [1.06274280 1.12849948]
% MC anti: est, std, conf. int: 1.105492851 0.003340705 [1.098945189 1.112040513]
% Relative reduction in standard deviation: 0.8008518
\begin{table}[!htb]
    \centering
    \begin{tabular}{|c|c|c|c|}
        \hline
        Estimator & Estimate & Standard Deviation & 95\% Confidence Interval \\
        \hline
        Standard Monte Carlo & 1.096 & 0.017 & [1.0627 1.128] \\
        Antithetic Monte Carlo & 1.106 & 0.0033 & [1.099 1.112] \\
        \hline
    \end{tabular}
    \caption{Rounded estimates for the integral $\int_{0}^{2} \frac{1}{1 + x^2} dx$ using the standard Monte Carlo estimator and the antithetic Monte Carlo estimator for $N = 100$ samples with seed $1234$}
    \label{tab:antithetic}
\end{table}

\begin{figure}[!htb]
    \centering
    \includesvg[width=1\textwidth]{figures/mc_anti.svg}\hfill
    \caption{Illustration of Monte Carlo Integration with Antithetic Variables. Notice the even distribution of points around $\frac{a+b}{2} = 1$ by construction yielding an advantage over the standard Monte Carlo estimator.}
    \label{fig:mc_antithetic}
\end{figure}

\subsubsubsection{Importance Sampling}

\problem{Consider we have a sharply peaked function $g(x)$ over the domain of integration. In this case using evaluation points sampled
from the uniform will be inefficient.}
\greenbox{\textbf{Importance sampling idea:} Preferably choose points where the integrand is large. This is done by sampling
from a distribution $f(x)$ similarly shaped as $g(x)$ but from which we can easily sample (e.g. by direct inversion\footnote{
Note that for the inverse transform we need the inverted CDF, so we need to integrate the PDF. So choosing $f(x)$ itself
as $g(x)$ is non-sense as we would need to integrate $f(x)$ to obtain the CDF.} or the rejection method).}

The higher the variance of the output of the function $g(x)$ over the domain of integration, the higher
the variance in our Monte Carlo Estimate for our integral. 

\begin{equation}
    \begin{multlined}
        I := \int_{V} g(x) dx = \int_{V} \frac{g(x)}{f(x)} f(x) dx, \quad I_N = \frac{1}{N} \sum_{i=1}^{N} \frac{g(x_i)}{f(x_i)} \\
        \text{sample} \left\{x_i\right\}_{i=1,...,N} \text{drawn from } f(x) \text{ limited to } V
    \end{multlined}
    \label{eq:importance_sampling}
\end{equation}

Let us as an example tackle the integral

\begin{equation}
    I = \int_{0}^{1} \frac{4}{1+x^2} dx \explain{=}{\text{exact solution}} \pi
\end{equation}

using the importance sampling function $f(x) = \frac{4 - 2x}{3}$ on $0 \leq x \leq 1$. We can sample
from $f(x)$ using the inverse transform method using the CDF $F(x) = \frac{4x - x^2}{3}$ (with $F(1) = 1$ as we want) so $F^{-1}(u) = 2 - \sqrt{4 - 3u}$. We can then
simply use equation \ref{eq:importance_sampling} to obtain an estimate.

The importance sampling approach compared to the standard Monte Carlo approach is illustrated in Fig. \ref{fig:mc_impA} and \ref{fig:mc_impB}. There it is easy to see how
effectively integrating a flatter function is to our advantage.

\greenbox{\textbf{Advantage of Importance Sampling}: As we can see in Fig. \ref{fig:mc_impA} the flatter $h(x) := \frac{g(x)}{f(x)}$ (best $h \propto g$ as achievable in lattice Monte Carlo) (mind the different naming in the figure) the sharper
the probability distribution $p(h)$ (a limited range of $h$ values with higher probabilities), the smaller the Monte Carlo error $\sigma_{\hat{I}_N} = \sqrt{\frac{\sigma_h^2}{N}}$.}

\begin{figure}[!htb]
    \centering
    \includesvg[width=1\textwidth]{figures/mc_imp_N_10.svg}\hfill
    \caption{Illustration of Monte Carlo Integration with Importance Sampling for $N=10$ samples.}
    \label{fig:mc_impA}
\end{figure}

\begin{figure}[!htb]
    \centering
    \includesvg[width=1\textwidth]{figures/mc_imp_N_200.svg}\hfill
    \caption{Illustration of Monte Carlo Integration with Importance Sampling for $N=200$ samples.}
    \label{fig:mc_impB}
\end{figure}

\pagebreak